%% kapa: exam=1 round=31 year=2020 period=1 form=A subject=civil_law subject_ko=민법 seq=11 number=11
%% source: Q-Net 감정평가사 기출문제 (cst00309) artlSeq=5207219 file=extracted/1008/5207219/A형/A형/2020년도 제31회 감정평가사 1차 시험_1교시_A형.pdf
%% conversion: kordoc 4.24.51 -> tools/gen1.py (auto); review=proofread-sonnet; flags=none
\begin{question}{11}
18세의 甲은 乙의 대리인을 사칭하여 그가 보관하던 乙의 노트북을 그 사정을 모르는 丙에게 팔았다. 이에 관한 설명으로 옳지 않은 것은? (다툼이 있으면 판례에 따름)
\begin{choices}
  \item 乙이 丙에게 매매계약을 추인한 때에는 매매계약은 확정적으로 효력이 생긴다.
  \item 乙이 甲에게 추인한 때에도 그 사실을 모르는 丙은 매매계약을 철회할 수 있다.
  \item 乙이 추인하지 않으면, 甲은 자신의 선택으로 丙에게 매매계약을 이행하거나 손해를 배상하여야 한다.
  \item 丙이 甲에게 대리권이 없음을 알았더라도 丙은 乙에게 추인 여부의 확답을 최고할 수 있다.
  \item 乙이 추인한 때에는 甲은 자신이 미성년자임을 이유로 매매계약을 취소하지 못한다.
\end{choices}
\end{question}
\answer{3}
